\documentclass[svgnames]{beamer}
\mode<presentation>
\usefonttheme{serif}
\usecolortheme{dove}
\useinnertheme{rounded}
%\useoutertheme{smoothbars}
\setbeamercolor{item projected}{fg=black}
\setbeamertemplate{navigation symbols}{}
%\setbeamertemplate{footline}[frame number] 

\usepackage[english]{babel}
\usepackage[latin1]{inputenc}
\usepackage{times}
\usepackage{amsthm,amssymb,amsmath,graphicx}
\usepackage{color}
\usepackage{gastex}

%\usepackage{pgfmath}
\usepackage{tikz,ifthen}%



%%%%%%%%%%%%%%%%%%%%%%%%%%%%%%%%%%%%%%%%%%%%%%%%%%%%%%%%%%%%%%%%%%%%%%%%%%%%%%%
%%%%%%%%%%%%%%%%%%%% A non-original creation by Nathana�l Fijalkow and myself %

\setbeamertemplate{frametitle}{%

  \vskip-2pt%

  \begin{beamercolorbox}[rightskip=2cm,leftskip=1em,dp=1ex,wd=12.8cm]{frametitle}%
    \vskip2pt%
    \usebeamercolor{frametitle}%
    \begin{tikzpicture}[scale=1]%
      \useasboundingbox (0,0) rectangle (0,0); %(-1,-1) rectangle (1,1);%
      \ifthenelse{\insertframenumber<\inserttotalframenumber}%
      { % uncomplete tart


        \pgfmathsetmacro{\aimangle}{90-(\insertframenumber*360/\inserttotalframenumber)}
        \fill [fill=frametitle.fg,thin, color=gray!50,draw=black] (11.8,.2) -- (11.8,.6) arc (90:\aimangle:0.4) -- cycle;%

      }{ % the full tart
        \fill[fill=frametitle.fg,thin, color=gray!50,draw=black] (11.8,0.2) circle (.4);%
      }%
      \fill[fill=frametitle.fg,thin, color=white,draw=black] (11.8,0.2) circle (.3);%
      \node at (11.8, .2) [black,circle]{\normalsize\insertframenumber};


    \end{tikzpicture}
    \insertframetitle%
    \vskip2pt%
  \end{beamercolorbox}%
}
%%%%%%%%%%%%%%%%%%%%%%%%%%%%%%%%%%%%%%%%%%%%%%%%%%%%%%%%%%%%%%%%%%%%%%%%%%%%%%%

%%% frametitle with tart-like progress meter
%\setbeamertemplate{frametitle}{%
%  \vskip-2pt%
%  \begin{beamercolorbox}[rightskip=2cm,leftskip=1em,dp=1ex,wd=12.8cm]{frametitle}%
%    \vskip2pt%
%    \usebeamercolor{frametitle}%
%    \begin{tikzpicture}[scale=1]% 
%      \useasboundingbox (0,0) rectangle (0,0); %(-1,-1) rectangle (1,1);%
%      \ifthenelse{\insertframenumber<\inserttotalframenumber}{%
%        %% uncomplete tart          
%        \pgfmathsetmacro{\aimangle}{90-(\insertframenumber*360/\inserttotalframenumber)}
%        \fill [fill=frametitle.fg,thin, color=gray!50] (11.5,.1) -- (11.5,.5) arc
%        (90:\aimangle:0.44) -- cycle;%
%        \fill [fill=frametitle.fg,thin, color=gray!25] (11.5,.1) -- (11.5,.5) arc
%        (90:\aimangle:0.38) -- cycle;%
%		\draw (11.5, .1) node[circle]{\begin{normalsize}\insertframenumber\end{normalsize}};
%      }{
%        %% the full tart
%        \fill[fill=frametitle.fg,thin, color=gray!50] (11.5,.053) circle (.44);%
%        \fill[fill=frametitle.fg,thin, color=gray!25] (11.5,.1) circle (.38);%
%		\draw (11.5, .1) node[circle]{\begin{normalsize}\insertframenumber\end{normalsize}};
%      }%
%    \end{tikzpicture}
%    \insertframetitle%
%    \vskip2pt%
%  \end{beamercolorbox}% 
%}

\renewcommand{\AA}{\mathcal{A}}

\newcommand{\PSPACE}{\mathrm{PSPACE}}
\newcommand{\BB}{\mathcal{B}}
\newcommand{\CC}{\mathcal{C}}
\newcommand{\JJ}{\mathcal{J}}
\newcommand{\DD}{\mathcal{D}}
\newcommand{\KK}{\mathcal{K}}
\newcommand{\LL}{\mathcal{L}}
\newcommand{\HH}{\mathcal{H}}
\newcommand{\GG}{\mathcal{G}}
\newcommand{\RR}{\mathbb{R}}
\newcommand{\NN}{\mathbb{N}}
\newcommand{\QQ}{\mathbb{Q}}
\newcommand{\MM}{\mathcal{M}}
\newcommand{\MMAA}{\MM_\AA}
\newcommand{\tr}[1]{\langle #1 \rangle}
\newcommand{\prob}[1]{\mathbb{P}_{#1}}
\newcommand{\stab}[1]{S(#1)}
\newcommand{\set}[1]{\{ #1 \}}
\newcommand{\val}[1]{\text{val}(#1)}
\newcommand{\merge}{\textrm{merge}}
\newcommand{\chck}{\textrm{check}}
\newcommand{\apply}{\textrm{apply}}
\newcommand{\finish}{\textrm{finish}}
\newcommand{\wait}{\textrm{wait}}
\newcommand{\delay}{\textrm{Delay}}
   
\title{Deciding the value $1$ problem \\ for probabilistic leaktight automata~\footnote{To appear in LICS~2012}
\vskip1em}
\author{Nathana\"el Fijalkow, \\
joint work with Hugo Gimbert and Youssouf Oualhadj}
\institute{LIAFA, Universit\'e Paris 7, France,\\
University of Warsaw, Poland,\\
\'Ecole Normale Sup\'erieure de Cachan, France.\\
}
\date{WATA, Dresden, Germany\\
May 31st 2012}

\setbeamertemplate{blocks}[rounded]%
%[shadow=true]
\setbeamercolor{block title}{bg=normal text.bg!90!black}
\setbeamercolor{block body}{bg=normal text.bg!95!black}

\AtBeginSection[]
{
  \begin{frame}{Outline}
    \tableofcontents[currentsection,currentsubsection]
  \end{frame}
\addtocounter{framenumber}{-1}
}
\AtBeginSubsection[]
{
  \begin{frame}{Outline}
    \tableofcontents[currentsection,currentsubsection]
  \end{frame}
\addtocounter{framenumber}{-1}
}

\begin{document}

\addtocounter{framenumber}{-1}
\begin{frame}[plain]
  \titlepage
%\begin{center}
%\includegraphics[scale=0.2]{warsaw.eps}
%\end{center}

\end{frame}

\begin{frame}<beamer>{Outline}
\addtocounter{framenumber}{-1}
  \tableofcontents
\end{frame}

\section{The value $1$ problem for probabilistic automata}

\begin{frame}{Probabilistic automata (Rabin, 1963)}
\begin{figure}
\begin{center}
\begin{picture}(20,40)(0,-10)
	\gasset{Nw=6,Nh=6}

  	\node[Nmarks=i,iangle=0](L1)(15,15){$1$}
  	\node(L2)(0,30){$2$}
  	\node[Nmarks=r](L3)(0,0){$3$}

  	\drawedge(L1,L2){$b$}
  	\drawedge[curvedepth=-5,ELside=r](L1,L3){$a,.4$}
  	\drawedge[curvedepth=-5,ELside=r](L3,L1){$b$}
	\drawloop(L1){$a,.6$}
	\drawloop[loopangle=135](L2){$a,b$}
	\drawloop[loopangle=215](L3){$a$}
\end{picture}
\end{center}
\end{figure}
$$\prob{\AA} : A^* \rightarrow [0,1]$$
\end{frame}

\begin{frame}{The value $1$ problem}
The value $1$ problem is, given a probabilistic automaton $\AA$:
\begin{center}
``are there words accepted by $\AA$ with arbitrarily high probability?''
\end{center}
\pause

\vskip2em
Define $\val{\AA} = \sup_w \prob{\AA}(w)$.
An equivalent formulation of this problem is:
\begin{center}
``$\val{\AA} \stackrel{?}{=} 1$''.
\end{center}

\pause
\begin{theorem}[Gimbert, Oualhadj, 2010]
The value $1$ problem is undecidable.
\end{theorem}
\end{frame}

\begin{frame}{An intuition}

\begin{figure}
\begin{center}
\begin{picture}(60,35)(0,0)
	\gasset{Nw=7,Nh=7}

  	\node[Nmarks=i,iangle=-90](0)(30,15){$0$}
  	\node(L1)(0,10){$L_1$}
  	\node[Nmarks=r](L2)(0,30){$L_2$}
  	\node(R1)(60,10){$R_1$}
  	\node(R2)(60,30){$R_2$}

	\drawloop(0){$a$}

  	\drawedge[curvedepth=5,ELside=l](0,L1){$b,\frac{1}{2}$}
  	\drawedge[curvedepth=5,ELside=l](L1,0){$a,1-x$}
  	\drawedge(L1,L2){$b$}
	\drawloop[loopangle=-135](L1){$a,x$}
	\drawloop[loopangle=90](L2){$a,b$}

  	\drawedge[curvedepth=-5,ELside=r](0,R1){$b,\frac{1}{2}$}
  	\drawedge[curvedepth=-5,ELside=r](R1,0){$a,x$}
  	\drawedge[ELside=r](R1,R2){$b$}
	\drawloop[loopangle=-45](R1){$a,1-x$}
	\drawloop(R2){$a,b$}
\end{picture}
\end{center}
\end{figure}
\pause
\begin{center}
has value $1$ if and only if $x > \frac{1}{2}$.
\end{center}
\end{frame}

%\begin{frame}{Our objective}
%Define a \textbf{large} and \textbf{interesting} subclass of probabilistic automata for which the value $1$ problem is decidable.
%\end{frame}

%\begin{frame}{A very restricted case}
%\begin{theorem}[Fijalkow, Gimbert, Oualhadj, 2011]
%The isolation problem is (still) undecidable if we randomise only on \textbf{one} transition.
%\end{theorem}
%\end{frame}

%\section{The limitedness problem for distance automata}
%
%\begin{frame}{Probabilistic automata VS distance automata}
%Consider a semiring $(\KK,+,\times)$.
%An automaton computes in the semiring $\KK$ if
%$\val{w} = \sum \set{\Pi(\rho) \mid \rho \textrm{ is a run over } w}$.
%
%\begin{itemize}
%	\item Probabilistic automata compute in $(\RR,+,\times)$ (there is a catch here).
%	\item Distance automata compute in the tropical semiring $(\NN \cup \set{\infty},\min,+)$.
%%	Here is an example:
%\end{itemize}
%
%%\begin{figure}
%%\begin{center}
%%\begin{picture}(40,15)(0,0)
%%	\gasset{Nw=6,Nh=6}
%%
%%  	\node[Nmarks=i,iangle=-90](p)(0,0){$p$}
%%  	\node[Nmarks=ir,iangle=-90](q)(20,0){$q$}
%%  	\node[Nmarks=r](r)(40,0){$r$}
%%
%%	\drawloop(p){$a,b:0$}
%%  	\drawedge(p,q){$b:0$}
%%	\drawloop(q){$a:1$}
%%  	\drawedge(q,r){$b:0$}
%%	\drawloop(r){$a,b:0$}
%%\end{picture}
%%\end{center}
%%\end{figure}
%\end{frame}
%
%\begin{frame}{The value $1$ problem VS the limitedness problem}
%The value $1$ problem for probabilistic automata is: 
%\begin{center}
%``are there words accepted with arbitrarily high probability?''.
%\only<1>{\phantom{undecidable}}
%\only<2>{\color{red}{undecidable}}
%\end{center}
%
%\vskip3em
%
%The \textbf{un}limitedness problem for distance automata is: 
%\begin{center}
%``are there words with arbitrarily high value?''.
%
%\pause
%\color{red}{decidable} 
%(Hashiguchi, 80's)
%\end{center}
%\end{frame}

\section{Towards an algebraic treatment of probabilistic automata}

%\begin{frame}{What kind of algebra?}
%Define an algebraic structure $\mathcal{M}$:
%\begin{figure}
%\begin{huge}
%\begin{center}
%\begin{picture}(60,40)(0,0)
%	\gasset{Nadjust=w}	
%
%	\node[Nframe=n](astar)(0,20){$A^*$}
%	\node[Nframe=n](monoid)(60,20){$\mathcal{M}$}
%
%	\drawedge[curvedepth=10](astar,monoid){$\textrm{abstraction}$}
%	\drawedge[curvedepth=10](monoid,astar){$\textrm{reification}$}
%\end{picture}
%\end{center}
%\end{huge}
%\end{figure}
%\pause
%What are the operations on $\mathcal{M}$?
%\end{frame}

\begin{frame}{Weighted automata using algebra (Sch\"utzenberger)}
\vspace*{1em}
\begin{figure}
\begin{center}
\begin{picture}(60,14)(0,0)
	\gasset{Nw=6,Nh=6,loopdiam=6}

  	\node[Nmarks=i,iangle=-90](0)(20,5){$0$}
  	\node(1)(40,5){$1$}
  	\node[Nmarks=r](F)(60,5){$F$}
  	\node(bottom)(0,5){$\perp$}

	\drawloop(0){$a,.3$}
	\drawloop(1){$a$}
	\drawloop[loopangle=0](F){$a,b$}
	\drawloop[loopangle=180](bottom){$a,b$}

  	\drawedge(0,bottom){$b$}
  	\drawedge[curvedepth=2](0,1){$a,.7$}
  	\drawedge[curvedepth=2](1,0){$b,.5$}
  	\drawedge(1,F){$b,.5$}
\end{picture}
\end{center}
\end{figure}
\vspace*{1.5em}

$$\tr{a} = 
\left(\begin{array}{cccc}
1 & 0 & 0 & 0 \\
0 & .3 & .7 & 0 \\
0 & 0 & 1 & 0 \\
0 & 0 & 0 & 1
\end{array}\right)
\qquad
\tr{b} = 
\left(\begin{array}{cccc}
1 & 0 & 0 & 0 \\
1 & 0 & 0 & 0 \\
0 & .5 & .5 & 0 \\
0 & 0 & 0 & 1
\end{array}\right)$$

\vspace*{1.5em}

\only<1>{
$$I = \left(\begin{array}{cccc} 0 & 1 & 0 & 0 \end{array}\right)
\qquad
F = \left(\begin{array}{c} 0 \\ 0 \\ 0 \\ 1 \end{array}\right)$$}

\only<2>{
$$I \cdot \tr{aaabaa} \cdot F = 
\prob{\AA}(aaabaa)$$
}
\vspace*{10em}
\end{frame}

\begin{frame}{Our objective}
Decide the value $1$ problem for a \textit{subclass} of probabilistic automata,
by \textbf{algebraic} and \textbf{non-numerical} means.
\pause

\begin{itemize}
	\item \textbf{algebraic:} focus on the automaton structure,
	\item \textbf{non-numerical:} abstract away the values.
\end{itemize}

\pause
Hence we consider non-deterministic automata:
we project $(\RR,+,\times)$ into the boolean semiring $(\set{0,1},+,\times)$.
\pause
\begin{figure}
\begin{center}
\begin{picture}(20,15)(0,5)
	\gasset{Nw=6,Nh=6}

  	\node[Nmarks=i,iangle=180](1)(0,5){$1$}
  	\node[Nmarks=r](2)(20,5){$2$}

	\drawloop(1){$a$}
  	\drawedge(1,2){$a$}
	\drawloop(2){$a$}
\end{picture}
\end{center}
\end{figure}
\end{frame}

\begin{frame}{Simplifying}
\vspace*{1em}
\begin{figure}
\begin{center}
\begin{picture}(60,14)(0,0)
	\gasset{Nw=6,Nh=6,loopdiam=6}

  	\node[Nmarks=i,iangle=-90](0)(20,5){$0$}
  	\node(1)(40,5){$1$}
  	\node[Nmarks=r](F)(60,5){$F$}
  	\node(bottom)(0,5){$\perp$}

	\drawloop(0){$a$}
	\drawloop(1){$a$}
	\drawloop[loopangle=0](F){$a,b$}
	\drawloop[loopangle=180](bottom){$a,b$}

  	\drawedge(0,bottom){$b$}
  	\drawedge[curvedepth=2](0,1){$a$}
  	\drawedge[curvedepth=2](1,0){$b$}
  	\drawedge(1,F){$b$}
\end{picture}
\end{center}
\end{figure}
\vspace*{1em}

$$\tr{a} = 
\left(\begin{array}{cccc}
1 & 0 & 0 & 0 \\
0 & 1 & 1 & 0 \\
0 & 0 & 1 & 0 \\
0 & 0 & 0 & 1
\end{array}\right)
\qquad
\tr{b} = 
\left(\begin{array}{cccc}
1 & 0 & 0 & 0 \\
1 & 0 & 0 & 0 \\
0 & 1 & 1 & 0 \\
0 & 0 & 0 & 1
\end{array}\right)$$

\vspace*{1em}

$$I \cdot \tr{u} \cdot F = 1 \quad \textrm{ if and only if } \quad \prob{\AA}(u) > 0$$
\end{frame}

\begin{frame}{The $\sharp$-expressions from Leung's algorithm}
\vspace*{3em}
\begin{figure}
\begin{center}
\begin{picture}(20,10)(0,0)
	\gasset{Nw=6,Nh=6}

  	\node[Nmarks=i,iangle=180](1)(0,5){$1$}
  	\node[Nmarks=r](2)(20,5){$2$}

	\drawloop(1){$a$}
  	\drawedge(1,2){$a$}
	\drawloop(2){$a$}
\end{picture}
\end{center}
\end{figure}
\vspace*{1.5em}
$$\only<1,2,3>{
\tr{a} = 
\left(\begin{array}{ccc}
1 & 1 \\
0 & 1
\end{array}\right)
	}
\qquad
\only<2,3>{
\tr{a^\sharp} = 
\left(\begin{array}{ccc}
0 & 1\\
0 & 1
\end{array}\right)}$$
\vspace*{1em}
In $\tr{a}$, the state $1$ is transient and the state $2$ is recurrent.

\vspace*{1em}
$$\only<3>{M^\sharp(s,t) = 
\left\{\begin{array}{ll}
1 & \textrm{if } M(s,t) = 1 \textrm{ and } t \textrm{ recurrent in } M,\\
0 & \textrm{otherwise.}
\end{array}\right.}$$
\vspace*{10em}
\end{frame}


\begin{frame}{The algorithm}
Compute a monoid inside the \textbf{finite} monoid $\MM_{Q \times Q}(\set{0,1},+,\times)$.

\begin{itemize}
	\item Compute $\tr{a}$ for $a \in A$:
$$\tr{a}(s,t) = 
\left\{\begin{array}{ll}
1 & \textrm{if } \prob{\AA}(s \xrightarrow{a} t) > 0,\\
0 & \textrm{otherwise.}
\end{array}\right.$$	
	\item Close under product and stabilization.
	\item \pause If there exists a matrix $M$ such that 
	$$\forall t \in Q, \quad M(s_0,t) = 1 \Rightarrow t \in F$$
	then ``$\AA$ has value $1$'',
	otherwise ``$\AA$ does not have value $1$''.
\end{itemize}
\end{frame}

\begin{frame}{An example}
\begin{figure}
\begin{center}
\begin{picture}(60,15)(0,-5)
	\gasset{Nw=6,Nh=6,loopdiam=6}

  	\node[Nmarks=i,iangle=-90](0)(20,5){$0$}
  	\node(1)(40,5){$1$}
  	\node[Nmarks=r](F)(60,5){$F$}
  	\node(bottom)(0,5){$\perp$}

	\drawloop(0){$a$}
	\drawloop(1){$a$}
	\drawloop[loopangle=0](F){$a,b$}
	\drawloop[loopangle=180](bottom){$a,b$}

  	\drawedge(0,bottom){$b$}
  	\drawedge[curvedepth=2](0,1){$a$}
  	\drawedge[curvedepth=2](1,0){$b$}
  	\drawedge(1,F){$b$}
\end{picture}
\invisible<1,5,6>{
\begin{picture}(60,10)(0,0)
	\gasset{Nw=5,Nh=5,loopdiam=4}

	\put(-20,4){$\tr{a}$}

  	\node(0)(20,5){$0$}
  	\node(1)(40,5){$1$}
  	\node[Nmarks=r](F)(60,5){$F$}
  	\node(bottom)(0,5){$\perp$}

	\drawloop(0){}
	\drawloop(1){}
	\drawloop[loopangle=0](F){}
	\drawloop[loopangle=180](bottom){}

  	\drawedge(0,1){}
\end{picture}}
\invisible<1,2,4,6>{
\begin{picture}(60,10)(0,0)
	\gasset{Nw=5,Nh=5,loopdiam=4}

	\put(-20,4){$\tr{b}$}

  	\node(0)(20,5){$0$}
  	\node(1)(40,5){$1$}
  	\node[Nmarks=r](F)(60,5){$F$}
  	\node(bottom)(0,5){$\perp$}

	\drawloop[loopangle=0](F){}
	\drawloop[loopangle=180](bottom){}

  	\drawedge(0,bottom){}
  	\drawedge(1,0){}
  	\drawedge(1,F){}
\end{picture}}
\invisible<1,2,3,6>{
\begin{picture}(60,10)(0,0)
	\gasset{Nw=5,Nh=5,loopdiam=4}

	\put(-20,4){$\tr{a^\sharp}$}

  	\node(0)(20,5){$0$}
  	\node(1)(40,5){$1$}
  	\node[Nmarks=r](F)(60,5){$F$}
  	\node(bottom)(0,5){$\perp$}

	\drawloop(1){}
	\drawloop[loopangle=0](F){}
	\drawloop[loopangle=180](bottom){}

  	\drawedge(0,1){}
\end{picture}}
\invisible<1,2,3,4>{
\begin{picture}(60,10)(0,0)
	\gasset{Nw=5,Nh=5,loopdiam=4}

	\put(-22,4){$\tr{a^\sharp \cdot b}$}

  	\node(0)(20,5){$0$}
  	\node(1)(40,5){$1$}
  	\node[Nmarks=r](F)(60,5){$F$}
  	\node(bottom)(0,5){$\perp$}

	\drawloop(0){}
	\drawloop[loopangle=0](F){}
	\drawloop[loopangle=180](bottom){}

  	\drawedge[curvedepth=-4](0,F){}
  	\drawedge(1,0){}
  	\drawedge(1,F){}
\end{picture}}
\invisible<1,2,3,4,5>{
\begin{picture}(60,10)(0,0)
	\gasset{Nw=5,Nh=5,loopdiam=4}

	\put(-25,4){$\tr{(a^\sharp \cdot b)^\sharp}$}

  	\node(0)(20,5){$0$}
  	\node(1)(40,5){$1$}
  	\node[Nmarks=r](F)(60,5){$F$}
  	\node(bottom)(0,5){$\perp$}

	\drawloop[loopangle=0](F){}
	\drawloop[loopangle=180](bottom){}

  	\drawedge[curvedepth=-4](0,F){}
  	\drawedge(1,F){}
\end{picture}}
\end{center}
\end{figure}
\pause\pause\pause\pause\pause\pause
\end{frame}

\begin{frame}{Correctness}
\begin{theorem}
If there exists a matrix $M$ such that 
$$\forall t \in Q, \quad M(s_0,t) = 1 \Rightarrow t \in F$$
then $\AA$ has value $1$.
\end{theorem}
\pause
But the value $1$ problem is undecidable, so\ldots
\end{frame}

\begin{frame}{No completeness}
\begin{figure}
\begin{center}
\begin{picture}(60,35)(0,0)
	\gasset{Nw=7,Nh=7}

  	\node[Nmarks=i,iangle=-90](0)(30,15){$0$}
  	\node(L1)(0,10){$L_1$}
  	\node[Nmarks=r](L2)(0,30){$L_2$}
  	\node(R1)(60,10){$R_1$}
  	\node(R2)(60,30){$R_2$}

	\drawloop(0){$a$}

  	\drawedge[curvedepth=5,ELside=l](0,L1){$b,\frac{1}{2}$}
  	\drawedge[curvedepth=5,ELside=l](L1,0){$a,1-x$}
  	\drawedge(L1,L2){$b$}
	\drawloop[loopangle=-135](L1){$a,x$}
	\drawloop[loopangle=90](L2){$a,b$}

  	\drawedge[curvedepth=-5,ELside=r](0,R1){$b,\frac{1}{2}$}
  	\drawedge[curvedepth=-5,ELside=r](R1,0){$a,x$}
  	\drawedge[ELside=r](R1,R2){$b$}
	\drawloop[loopangle=-45](R1){$a,1-x$}
	\drawloop(R2){$a,b$}
\end{picture}
\end{center}
\end{figure}

%For all matrices $M$,
%$$\exists t \in Q, \quad M(s_0,t) = 1 \wedge t \notin F$$
%but $\AA$ has value $1$. 

Left and right parts are symmetric, so for all $M$:
$$M(0,L_2) = 1 \Longleftrightarrow M(0,R_2) = 1.$$
\end{frame}

\begin{frame}{A leak}
\begin{figure}
\begin{center}
\begin{picture}(50,30)(0,0)
	\gasset{Nw=6,Nh=6}

  	\node[Nmarks=i,iangle=0](L1)(15,15){$1$}
  	\node(L2)(0,30){$2$}
  	\node[Nmarks=r](L3)(0,0){$3$}

  	\drawedge(L1,L2){$b$}
  	\drawedge[curvedepth=-5,ELside=r](L1,L3){$a$}
  	\drawedge[curvedepth=-5,ELside=r](L3,L1){$b$}
	\drawloop(L1){$a$}
	\drawloop[loopangle=135](L2){$a,b$}
	\drawloop[loopangle=215](L3){$a$}

	\drawline[linewidth=.3,AHnb=0](26,-3)(26,40)

\pause

	\put(45,-10){$\tr{a^\sharp \cdot b}$}

  	\node(L1b)(55,15){$1$}
  	\node(L2b)(40,30){$2$}
  	\node[Nmarks=r](L3b)(40,0){$3$}

  	\drawedge(L3b,L1b){}
	\drawloop(L1b){}
	\drawloop[loopangle=135](L2b){}

\only<3>{\drawedge[dash={1}0](L1b,L2b){$\varepsilon$}}
\end{picture}
\end{center}
\end{figure}
\pause
\vskip2em
\begin{center}
There is a leak from $1$ to $2$.
\end{center}
\pause
\begin{definition}
An automaton $\AA$ is leaktight if it has no leak.
\end{definition}
\end{frame}

\begin{frame}{Leaktight automata}

\begin{theorem}[Fijalkow,Gimbert,Oualhadj]
The value $1$ problem is decidable for leaktight automata.
\end{theorem}

\vskip2em
The proof relies on Simon's factorization forest theorem.
\end{frame}

\begin{frame}{Conclusion and perspectives}
\vskip-1em
\begin{itemize}
	\item We defined a subclass of probabilistic automata which subsumes all subclasses of probabilistic automata
whose value $1$ problem is known to be decidable,\vskip2em\pause
	\item We defined an algebraic algorithm (similar to Leung's algorithm) for the value $1$ problem 
and proved its completeness for the class of leaktight automata.\vskip2em\pause
	\item What does this algorithm actually compute? \vskip2em\pause
	\item Can we use similar algorithms for other semirings?
\end{itemize}
\end{frame}

\end{document}